\scanpage{6}
a) $\mathcal M_n=$ parties de cardinal $\le1$ de $R^n$

b) $\mathcal M_n=$ parties finies de $R^n$

c) $\mathcal M_n=$ toutes les parties de $R^n$

\margin{ici les axiomes a), b), c) du NB \textcircled{2}, \uline{à l'exclusion} de $c_2)$}

Aucun de ces exemples n'est particulièrement \unsure{passionnant}.

\margin{\textcircled{II} \uline{espaces modérés}\par
Dans ce qui suit \unsure{au lieu de} $(M1)$ il suffit l'axiome concernant \deleted{\unsure{injections|les applications}} \inserted{\unsure{les surjections}} \uline{pour les seuls} $[1,n]$ -- \unsure{c.-à-d. les} \unsure{linéaires des s.} \illegible{3} \unsure{applications} \unsure{des} \uuline{\unsure{directes}}.}
\margin{on n'exclut \unsure{pas} $I$ vide}

Terminologie : partie modérée de $R^I$, points modérés de $R$, de $R^I$.

\uline{déf} Soit $A\in\mathcal M_I$, $B\in\mathcal M_J$, $f:A\unsure{\longrightarrow|-}B$, d'où
\[
\Gamma_f\subset A\times B\subset R^I\times R^J\simeq R^{I\unsure{\cup|\amalg}J}.
\]
On dit que $f$ est \uline{modérée} si
\[
\Gamma_f\in\mathcal M_{I\unsure{\cup|\amalg}J}.
\]

\margin{\unsure{via|vu} $M2$}
\uline{Proposition} a) $\mathrm{id}_A$ (si $A\in\mathcal M_I$) est modérée

b) \diagram{\begin{tikzcd}[row sep=tiny,column sep=large]
f:A\arrow[r,no head]&B&\text{et}&g:B\arrow[r,no head]&C\\
\mathcal M_I\arrow[u,phantom,"\rotatebox{90}{\(\in\)}" description]&\mathcal M_J\arrow[u,phantom,"\rotatebox{90}{\(\in\)}" description]&&&\mathcal M_{\unsure{K|n}}\arrow[u,phantom,"\rotatebox{90}{\(\in\)}" description]
\end{tikzcd}} modérés $\Longrightarrow gf$ modéré

c) \diagram{\begin{tikzcd}[row sep=tiny]
f:A\arrow[r]&B\\
\mathcal M_I\arrow[u,phantom,"\rotatebox{90}{\(\in\)}" description]&\mathcal M_J\arrow[u,phantom,"\rotatebox{90}{\(\in\)}" description]
\end{tikzcd}} modéré et bijectif $\Longrightarrow f^{-1}$ modéré

D'où catégorie des \uline{modèles modérés} $\mathrm{Mod}_{\mathcal M_*}$, \inserted{(on prend $(I,A\in\mathcal M_I)$)} et un foncteur \uline{conservatif} \uline{fidèle}
\diagram{\begin{tikzcd}
\mathrm{Mod}_{\mathcal M_*}\arrow[r]&(\mathrm{Ens})
\end{tikzcd}}\qquad$(I,A)\longmapsto A$.
\margin{on n'exclut pas les modèles $\subset R^0=\{e\}$, \inserted{modèle} \deleted{\unsure{espace|discret}} réduit à un point \ldots}

On en déduit\deleted{\unsure{on}} la catégorie des ensembles \inserted{(ou ``espaces'')} $(\mathcal M_*-)$ modérés, équivalente à : $\mathrm{Mod}_{\mathcal M_*}$, mais avec un foncteur
\diagram{\begin{tikzcd}
\deleted{\unsure{\mathrm{Mod}}}\ \mathrm{Esp}_{\mathcal M_*}\arrow[r]&(\mathrm{Ens})
\end{tikzcd}}\qquad$X\longmapsto|X|$
qui est fidèle, conservatif et de plus ``transportable'', i.e.
\mspage{5}
